% !Mode:: "TeX:UTF-8"

\section{课题来源及研究的目的和意义}

\subsection{课题来源或研究背景}

随着大语言模型（LLM）应用边界的扩展，基于 LLM 的城市旅游行程规划系统
成为文旅智能化的典型场景：用户以自然语言表达天数、预算、偏好等需求，
系统输出逐日景点（POI）路线。然而，此类系统在真实落地中暴露出三个
尚未解决的问题。其一，\textbf{缺乏真实语料接地导致的幻觉}：LLM 会编造
不存在的地点与口碑、时序等事实细节；检索增强生成（RAG）是主流对策，
但旅游场景下"真实游记经验与 POI 属性的文本级对齐语料"本身缺失，
可检索的证据基础薄弱。其二，\textbf{空间行为接地不足}：在"从候选景点
中选择下一站"这类空间决策上，微调后的开源 LLM 的选择质量并不稳定优于
"总选最近的候选"这一简单规则——通用预训练不能替代对真实人群移动
统计规律的学习，需要以后训练注入领域行为。其三，\textbf{评估方法学
薄弱}：LLM 行程规划的评估普遍缺少与强规则基线的配对统计检验，对采样
方差与聚合谬误缺乏控制，概率校准几乎无人检验，"LLM 是否真的更好"
的结论易被评估缺陷污染。

本课题依托一个已运行的哈尔滨旅游行程规划系统（Qwen3.5-4B 微调 +
检索增强 + 规则层），将其升级为以"检索增强与后训练"为主线的规划
方法研究对象。系统研制过程中积累了两项数据资产：经语义标注的
10{,}000 个城市 POI 库（含空间聚类与距离矩阵），以及 168 条真实游记
路线（共 1{,}660 次相邻景点转移）——后者从未参与任何模型训练，是本
课题的行为学金标准。前期分析发现真实游客的转移距离呈显著双峰形态：
约半数转移发生在 2km 内的商圈尺度，其余为跨区长距离移动（中位数
1.64km，最大 195km）。围绕该数据资产，前期已完成可解释行为先验的
统计建模（二阶区域转移与混合距离衰减）与三层统计评测体系——在
新框架中，它们分别转化为"统计重排融合层"与"统计评测与校准层"的
前期基础：前者在前期实验中使融合命中率相对最强规则基线提升约
11 个百分点，后者为全部方法提供预注册的统一裁决。

\subsection{研究的目的及意义}

本课题的总体目标：以真实旅游行为数据为金标准，建立"检索增强
$\times$ 后训练 $\times$ 统计重排与校准"的大语言模型行程规划方法及其
三层统计评测体系。具体目标为：
\begin{enumerate}
  \item 构建行为接地的 RAG 语料与混合检索机制：将真实游记经验
        （小红书笔记的自由文本）对齐到 POI 库形成可检索的证据语料，
        抑制幻觉并向生成过程注入真实口碑与空间结构信息；
  \item 建立面向行程规划的 LLM 后训练方法：行程级监督微调（QLoRA）
        与以行为对数似然构造偏好对的直接偏好优化（DPO），并以统计
        先验重排融合约束模型输出的空间行为；
  \item 建立三层统计评测体系（点层配对检验、个体层似然、群体层分布
        距离），克服既有评价指标与就近基线同源等方法学缺陷。
\end{enumerate}

理论意义：其一，真实游记到 POI 的证据对齐给出旅游领域 RAG 语料的
构建与质量评估方法，是检索增强在垂直场景落地的关键一环；其二，以
统计似然构造偏好标签的后训练目标，将行为统计注入大模型优化本身，
为"统计方法与大模型结合"提供新的实证形态；其三，目的地内旅游转移
距离的双峰性在旅游移动文献中缺少正式检验报道，距离衰减"指数还是
幂律"的长期争论将在行程内转移粒度获得数据驱动裁决。应用意义：
框架直接服务于行程规划系统的工程部署；三层评测方法学可迁移至其他
"LLM 对比规则"的决策场景。

\section{国内外在该方向的研究现状及分析}

\subsection{国外研究现状}

\textbf{（1）下一站 POI 推荐。}第一代方法以马尔可夫链与矩阵分解为
主：因子化个性化马尔可夫链 FPMC\cite{rendle2010fpmc}、其空间扩展
FPMC-LR\cite{cheng2013fpmclr} 与个性化排名度量嵌入
PRME\cite{feng2015prme} 参数省、可解释，但一阶马尔可夫假设很少被
正式检验。第二代深度序列模型以 ST-RNN\cite{liu2016strnn}、
ARNN\cite{liu2019arnn}、长短短期偏好模型\cite{li2022lstpm}为代表，
在签到基准上精度领先，但数据饥渴、黑箱。第三代 LLM 方法将推荐转为
问答或生成任务\cite{li2024llmpoi,wu2024llm4rec}，近期工作进一步以
多因子协作引导 LLM 完成下一站推理\cite{song2026mfc4poi}，零样本
泛化与语义理解能力强，但存在四类系统性不足：幻觉编造
POI、空间结构弱（本项目前期实测：纯 LLM 在空间就近性上系统性输给
规则）、评估缺少与强规则基线的配对统计检验、概率校准被
忽视\cite{guo2017calibration}。

\textbf{（2）检索增强生成与 LLM 行程规划。}检索增强生成通过外挂
知识库为 LLM 提供事实接地\cite{lewis2020rag}，已成为抑制幻觉的主
流框架。旅游方向上，TravelRAG\cite{song2024travelrag} 以多层知识
图谱组织景点检索；STrajRAG\cite{lin2025strajrag} 用监督式轨迹检索
增强下一 POI 生成，与本课题场景最接近；TourMLM\cite{yamanishi2025tourmlm}
将检索增强扩展到多模态多任务。评估侧出现时空细粒度基准
TripCraft\cite{chaudhuri2025tripcraft}、个性化真实基准
TripTailor\cite{wang2025triptailor} 与行程校验护栏
Iti-Validator\cite{gadbail2025validator}。系统侧以"LLM 出初稿 +
算法修正"的混合架构为主流，如 RETAIL\cite{deng2025retail} 将 LLM
与蚁群优化结合，Google 亦采用混合行程规划架构\cite{google2025trip}。
共性不足：检索语料以结构化属性为主，真实游记经验与 POI 库的文本级
对齐缺失；评估以端到端打分为主，缺少"每一站是否符合真实人类流动
统计"的逐点检验。

\textbf{（3）大模型后训练。}参数高效微调以 LoRA\cite{hu2022lora}
与低比特量化的 QLoRA\cite{dettmers2023qlora} 为主流，使单卡低资源
适配成为可能；偏好优化方面，直接偏好优化 DPO\cite{rafailov2023dpo}
绕开奖励模型直接以偏好对训练。但在行程规划任务上，现有后训练多为
点级候选排序的指令微调，缺少行程级序列监督，更缺少以真实行为统计
为信号构造偏好的工作——后训练目标与"行为符合度"之间缺乏桥梁。

\textbf{（4）本文统计层的两个方法学支柱。}其一，概率融合与校准：
Product of Experts\cite{hinton2002poe} 与对数意见池化理论为"乘性
融合两个概率模型"提供经典框架，FuseLLM\cite{wan2024fusellm} 将其
用于 token 分布层融合；多专家学习延迟\cite{mao2023defer}与 LLM
路由\cite{ong2024routellm}研究"何时信任哪个模型"；对齐后的 LLM
常过度自信\cite{guo2017calibration}——这些构成本课题"统计先验
重排融合与校准层"的背景。其二，人类移动的距离衰减与分布形态：
幂律与指数之争由来已久\cite{chen2015decay}，个体移动距离重尾已有
经典刻画\cite{gonzalez2008understanding,simini2012universal}，旅游
流距离衰减见于旅游研究\cite{mckercher2003distance}——但目的地内
旅游转移距离的分布形态（尤其双峰性）缺少正式检验。

\subsection{国内研究现状}

国内团队在 LBSN 下一站推荐方向有系统性贡献：FPMC-LR\cite{cheng2013fpmclr}
与 PRME\cite{feng2015prme}均为国内团队的代表性工作。LLM 方向，
国内进展活跃：面向真实约束的 RETAIL 框架\cite{deng2025retail}将
LLM 与蚁群优化结合；《智能系统学报》近期发表了 LLM 推荐系统
综述\cite{xie2025llmrec}与 LLM 个性化推荐研究\cite{wu2024llmpersorec}，
系统梳理了提示工程、模型微调与范式转变。总体上，现有工作以系统
研制与综述为主，"行为接地的检索语料构建、面向行程规划的后训练、
以及概率级统计重排与严格评测"的研究较少。

\subsection{国内外文献综述的简析}

综合现状可得三个缺口：其一，行为接地的检索语料缺失——游记富文本
与 POI 库没有对齐，RAG 证据无法支撑行为合理的规划，真实转移的
空间结构（如重尾）难以进入检索层；其二，面向行程规划的后训练
不足——行程级序列监督与"以真实行为统计构造偏好"的后训练目标
缺位；其三，评估缺数据纪律——配对统计检验、聚合谬误与过拟合控制、
概率校准检验在 LLM 行程规划评估中普遍缺位。本课题的 RAG 语料
构建、似然驱动偏好优化与三层评测体系分别针对这三个缺口展开。

\section{前期的理论研究与试验论证工作的结果}

前期已完成核心数据资产建设、行为先验建模、LLM 系统与后训练管线、
评测体系设计与全部前置实验（代码与结果可复现）：
\begin{enumerate}
  \item \textbf{数据资产与双峰性检验}：168 条真实游记路线分层三分
        （拟合 84 / 权重学习 42 / 测试 42），三份分布一致性通过 KS
        检验；配对 McNemar 功效分析确定测试集需覆盖全部约 372 个
        预测点，最小可检差异约 6 个百分点。Hartigan dip
        test\cite{hartigan1985dip}在 $p<10^{-4}$ 水平拒绝单峰；
        跨区分量需重尾族，最终采用 lognormal 与 log-gamma 的异质
        混合；
  \item \textbf{行为先验建模（现统计重排层的构件）}：参数自助法
        似然比检验支持二阶区域转移（$p=0.005$，样本外对数损失
        0.62 对 0.75 nat/转移）；候选池离散选择任务上的条件 logit
        拟合显示混合距离衰减形式优于任一单一形式，工程常用的纯
        指数形式被数据否决；
  \item \textbf{LLM 系统与后训练管线（新主线工程底座）}：Qwen3.5-4B
        4-bit 量化 + QLoRA 监督微调的完整训练与推理链已在单卡
        RTX 4090 打通（一轮约 1 小时），编号候选 prompt 与首 token
        概率接口实现确定性候选打分，DPO/GRPO 变体亦已验证可训练；
  \item \textbf{统计重排融合的前期证据}：测试集 372 点 8 选 1 任务
        上，统计先验单独命中率 37.1\%（对就近规则 32.3\%，
        $p=10^{-4}$）；LLM 单独仅 26.6\%；两者 log-linear 池化后达
        43.3\%（$p\approx0$）——LLM 的价值在统计重排框架下实现
        反转。架构消融显示距离衰减为命脉组件（去除后掉 18.8 个
        百分点）；
  \item \textbf{三层评测体系 v1.0}：以构念效度理论区分"行为符合度"
        与"工程质量"，预注册三层裁决规则（点层 McNemar 主裁决、
        个体层似然描述性、群体层分布距离），并完成生成层首轮流证
        （融合方法群体分布距离最优，KS $=0.219$）；
  \item \textbf{负结果}：状态依赖融合权重与固定权重无显著差异
        （$p=1.0$），学得的权重塌缩为窄带常数；LLM 概率温度校准
        亦无实质增益（期望校准误差原始仅 0.041）。两者共同说明
        融合收益来自两模型误差互补的结构本身。
\end{enumerate}

\section{学位论文的主要研究内容、实施方案及其可行性论证}

\subsection{主要研究内容}

论文将围绕"检索增强语料—LLM 后训练—统计重排与校准—评测体系"
主线组织四章内容（研究内容以论证分析为主，非目录罗列）：

\begin{enumerate}
  \item \textbf{行为接地的 RAG 语料构建与混合检索}：POI 结构属性的
        卡片化文本合成；小红书游记证据到 POI 的对齐（别名归一、
        模糊匹配与区域约束，保守阈值加人工抽样校验并报告语料
        质量准确率）；稠密向量召回与规则打分的两路混合检索；检索
        质量评测（证据覆盖率、候选命中排名）与"纯规则/纯稠密/混合"
        消融。
  \item \textbf{面向行程规划的 LLM 后训练}：带证据 prompt 的行程级
        监督微调数据重建与 QLoRA 训练；以行为对数似然高低构造
        路线级偏好对的直接偏好优化（统计似然$\to$偏好标签$\to$
        后训练目标）；以后训练方法在三层评测与融合增益上的表现
        为裁决的评估协议。
  \item \textbf{统计先验重排融合与概率校准}：log-linear 池化的权重
        语义与性质；固定权重、状态依赖权重（逻辑回归堆叠、嵌套
        交叉验证）与硬路由的统一对比；LLM 候选概率的校准性分析
        （温度缩放、期望校准误差）；重排层组件必要性消融（距离
        衰减、区域粒度、转移阶数逐一去除检验）。
  \item \textbf{三层统计评测与系统集成验证}：点层（逐站命中率、
        Wilson 区间\cite{wilson1927probable}、配对精确
        McNemar\cite{mcnemar1947note}、多重比较
        校正\cite{holm1979simple}）；个体层（路线行为对数似然，
        配对 Wilcoxon，标注其模型依赖性）；群体层（转移距离与类型
        节奏的 KS/Wasserstein 距离\cite{massey1951ks}、路线级自助
        置信区间）；以构念效度理论\cite{cronbach1955construct}区分
        "行为符合度"与"工程质量"两类指标，三层互证以控制循环
        论证风险——该思路借鉴生成模型评测中似然与样本分布判据
        不可互替的论证\cite{theis2016note}。将 RAG、后训练与重排
        选择器接入行程规划管线，在生成层复验三层指标；案例分析
        与消融展示。
\end{enumerate}

\subsection{实施方案及其可行性论证}

技术路线已验证：后训练与推理全链在单卡 RTX 4090 完成（LLM 4-bit
量化推理，候选打分单次前向）；语料构建为纯数据工程，无算力风险；
统计重排的增益已在逐站层与生成层获得一致的前期证据（43.3\% 对
32.3\%，群体分布距离最优）。数据资产与代码均已版本化。可行性论证
同时包含对方法本身的辩证分析：

\textbf{方法合理性的有利证据}：重排融合增益在点层与生成层方向一致；
在已知偏向就近基线的旧工程指标下先验系方法仍然胜出，说明结论不
依赖单一指标族；RAG 与后训练的每个增量均可通过三层评测独立裁决。

\textbf{方法局限与对策}：（1）个体似然层与重排方法共用行为统计，
存在循环论证风险——对策是三层互证，裁决只依赖两层无模型指标；
（2）168 条真实路线是样本硬上限，效应量检验功效受限——对策是
功效前置设计与路线级自助；（3）游记证据对齐存在误配风险，错误
证据会向 LLM 注入噪声——对策是保守匹配阈值、人工抽样校验并在
消融中量化语料质量的影响；（4）全部生成方法都压缩了真实转移的
重尾（生成分位 90 约 5km 对真实约 8km），源于候选池以当前中心
检索——混合检索的全城召回本身就是针对性的缓解实验；（5）状态
权重与校准的负结果提示融合收益的上限来自误差互补结构——论文
将其定位为机理解释而非缺陷。

\section{论文进度安排，预期达到的目标}

\subsection{进度安排}

\begin{enumerate}
  \item 2026.09--2026.10：完成开题；混合模型参数区间估计与分量数
        选择收尾；游记证据到 POI 的对齐与 RAG 语料构建；检索质量
        评测。
  \item 2026.11--2026.12：混合检索接入；带证据 prompt 的行程级
        SFT 与重排融合复验。
  \item 2027.01--2027.02：中期检查；行为似然偏好对的 DPO 实验；
        全部消融复跑与多重比较校正固化。
  \item 2027.03--2027.04：论文撰写与图表规范化；预答辩。
  \item 2027.05--2027.06：修改完善、查重、答辩。
\end{enumerate}

\subsection{预期达到的目标}

\begin{enumerate}
  \item 建立目的地内旅游转移距离双峰性的正式建模与检验方法；
  \item "RAG + 后训练 + 统计重排"框架在真实行为数据上显著超越
        LLM 基线，且不劣于已获融合结果的前期水平（43.3\%）；
  \item 形成可复用的 LLM 空间推荐统计评测方法学与开源代码；完成
        学位论文一篇，投稿学术论文一篇。
\end{enumerate}

\section{学位论文预期创新点}

\begin{enumerate}
  \item \textbf{实证层面}：首次对目的地内旅游转移距离给出双峰性的
        正式检验（$p<10^{-4}$）与重尾混合参数化，并在行程粒度对
        距离衰减函数形式给出数据驱动裁决。
  \item \textbf{方法层面}：提出行为接地的 RAG 语料构建方法（真实
        游记证据到 POI 的对齐与质量声明）与统计似然驱动的 DPO 偏好
        构造，形成"检索增强—后训练—统计重排"协同的 LLM 行程规划
        框架；实证"LLM 单独弱于规则、经统计重排后显著反超"的价值
        反转，并通过组件消融论证最终架构的最简形式。
  \item \textbf{评测层面}：构建三层统计评测体系，系统刻画 LLM 对比
        规则评估中的同源性、聚合谬误与过拟合陷阱及其对策，以构念
        效度理论明确"行为符合度"与"工程质量"的边界。
\end{enumerate}

\section{为完成课题已具备和所需的条件、外协计划及经费}

已具备：168 条真实游记路线与万级 POI 数据资产；全部实验代码与
结果（版本管理）；RTX 4090 算力；基座模型与微调适配器（本地推理，
无 API 开销）；中文核心文献已检索补齐。所需：可能的扩展数据用于
外部验证。经费需求为常规打印与差旅。

\section{预计研究过程中可能遇到的困难、问题，以及解决的途径}

\begin{enumerate}
  \item \textbf{语料对齐精度}：游记自由文本与 POI 的名称匹配存在
        误配与漏配，错误证据会污染检索与微调。对策：保守阈值与
        人工抽样校验（报告准确率）；消融中量化语料质量影响；
        若对齐覆盖率不足，降级为"POI 卡片 + 拟合份路线模式"语料。
  \item \textbf{样本量上限}：真实路线仅 168 条，细粒度检验功效受限。
        对策：功效前置设计；路线级自助；必要时以半参数模拟补充。
  \item \textbf{循环论证风险}：个体似然与重排方法共用行为统计。对策：
        三层互证，裁决依赖无模型层；似然层定位为描述性。
  \item \textbf{重尾压缩}：生成路线缺失真实数据中的极端长转移。对策：
        如实报告为共同局限；混合检索的全城召回作为针对性缓解
        实验并单独评测。
  \item \textbf{微调的两难}：在含距离线索的证据与候选上微调，可能
        使 LLM 模仿就近启发式而损失语义互补性。对策：以融合增益
        （而非 LLM 单独精度）为后训练的裁决指标。
\end{enumerate}

\section{主要参考文献}

\bibliographystyle{hithesis}
\bibliography{reference}
